\documentclass[12pt,a4paper]{amsart}

\usepackage[utf8]{inputenc}
\usepackage[T1]{fontenc}
\usepackage{amssymb,amsmath,amsthm}
\usepackage{hyperref}
\usepackage{geometry}
\usepackage{mathrsfs}
\usepackage{stmaryrd}

\newtheorem{theorem}{Theorem}[section]
\newtheorem{lemma}[theorem]{Lemma}
\newtheorem{corollary}[theorem]{Corollary}
\newtheorem{proposition}[theorem]{Proposition}
\theoremstyle{definition}
\newtheorem{definition}[theorem]{Definition}
\newtheorem{example}[theorem]{Example}
\theoremstyle{remark}
\newtheorem{remark}[theorem]{Remark}

\newcommand{\cat}{\mathcal}
\newcommand{\Set}{\mathbf{Set}}
\newcommand{\Preord}{\mathbf{Preord}}
\newcommand{\Cat}{\mathbf{Cat}}
\newcommand{\CAT}{\mathbf{CAT}}
\newcommand{\op}{\mathrm{op}}
\newcommand{\Ob}{\operatorname{Ob}}
\newcommand{\Hom}{\operatorname{Hom}}
\newcommand{\Nat}{\operatorname{Nat}}
\newcommand{\Id}{\operatorname{Id}}
\newcommand{\Tr}{\operatorname{Tr}}
\newcommand{\Aut}{\operatorname{Aut}}
\newcommand{\Iso}{\operatorname{Iso}}
\newcommand{\Eqv}{\operatorname{Equiv}}

\title{Exact enumeration of finite categories}
\author{Ponensmoder}
\date{2026}

\begin{document}

\maketitle

\begin{abstract}
  We present an algorithm for the exact enumeration of all finite categories
  up to a given total number of morphisms.  The algorithm stratifies
  enumeration over three enrichment layers --- the two-element lattice
  $\{0,1\}$, the natural numbers $\mathbb{N}$ with $\min$, and the category
  $\Set$ --- and exploits a decomposition theorem for categories enriched
  over a preorder.  We prove correctness via a colimit--limit duality for
  the functor $U(-)\text{-Cat} : \Preord^{\op} \to \CAT/\Set$, and we
  report verified counts up to 12 morphisms, including Cauchy-complete
  categories under strict isomorphism, categorical equivalence, labelled-object
  equality, and source--target duality.
\end{abstract}

\tableofcontents

\section{Introduction}
\label{sec:introduction}

The enumeration of finite algebraic structures is a classical problem in
computational mathematics.  While finite groups \cite{besche2002},
semigroups \cite{distler2011}, and lattices \cite{jevons1897} have received
extensive attention, the enumeration of finite categories --- the simplest
algebraic structure combining composition, identity, and associativity ---
has progressed more slowly.  Cruttwell \cite{cruttwell2014} gave the first
systematic counts up to 7 morphisms using a backtracking search over
composition tables, and the SmallCategories project \cite{smallcats}
extended these results with a database of canonical forms.

This paper presents a new algorithm for enumerating finite categories by
total morphism count, with three distinguishing features:

\begin{enumerate}
\item \textbf{Stratified enrichment.}  Rather than searching the space of
  composition tables directly, we decompose the problem into three layers:
  enrichment over the lattice $\{0,1\}$ (support preorders), enrichment
  over $\mathbb{N}$ with $\min$ (hom-count matrices), and enrichment over
  $\Set$ (composition tables).  Each layer is enumerated independently and
  refines the previous one.
\item \textbf{Decomposition-based generation.}  Categories are decomposed
  into \emph{biconnected components} (categories where every hom-set is
  nonempty) connected by profunctor edges.  This reduces the search space
  dramatically for larger morphism counts.
\item \textbf{Multiple quotient relations.}  We support four distinct
  notions of equivalence: strict isomorphism, categorical equivalence,
  labelled-object equality (where object labels matter but morphism labels
  within hom-sets do not), and source--target duality (identifying a
  category with its opposite).
\end{enumerate}

All code implementing the algorithms described here is publicly available
at \url{https://github.com/Leucocholia/cat-enum}.

\subsection{Related work}

Cruttwell \cite{cruttwell2014} enumerated finite categories up to 7
morphisms using a direct shape-based search: generate all hom-count
matrices satisfying support-transitivity, then fill the composition table
via backtracking with associativity propagation.  Our shape-based search
(Section~\ref{sec:shape-search}) follows this approach but adds
Cauchy-completeness pruning during search rather than after.

The SmallCategories project \cite{smallcats} maintains a database of
canonical keys for finite categories up to 9 morphisms.  Our
canonicalization algorithm (Section~\ref{sec:canonical}) follows the
Cruttwell--Leblanc ``arrows-only'' representation and exhaustive
relabelling under object and morphism permutations.

Baez and Dolan \cite{baez2001} observed that finite categories can be
viewed as categories enriched over the monoidal poset $(\mathbb{N}, \min,
+\infty)$, but did not pursue enumeration.  Our stratification into three
enrichment layers makes this observation algorithmic.

\subsection{Notation and conventions}

We work with a left-to-right composition convention: $\Hom(y,z) \times
\Hom(x,y) \to \Hom(x,z)$ sends $(g,f)$ to $f;g$.  All categories in this
paper are small and finite.  A \emph{finite category} $\cat{C}$ is
determined by a finite set $\Ob(\cat{C})$ of objects and a finite set
$\Mor(\cat{C})$ of morphisms, together with source and target functions
$s,t : \Mor(\cat{C}) \to \Ob(\cat{C})$, an identity assignment
$\id : \Ob(\cat{C}) \to \Mor(\cat{C})$, and a partial composition
operation $; : \Mor(\cat{C}) \times \Mor(\cat{C}) \rightharpoonup
\Mor(\cat{C})$ satisfying the usual axioms.

The \emph{total morphism count} $n = |\Mor(\cat{C})|$ is our primary size
parameter.  We enumerate categories by increasing $n$, and for each $n$ we
break down results by object count $k = |\Ob(\cat{C})|$.

\section{Enrichment lattices and the three-layer stratification}
\label{sec:enrichment}

The central observation of this paper is that a finite category can be
built in three successive enrichment steps, each over a simpler base.

\begin{definition}
  A \emph{monoidal poset} is a poset $(L, \leq)$ together with a binary
  operation $\wedge : L \times L \to L$ (the \emph{meet}) and a distinguished
  element $\top \in L$ (the \emph{unit}) such that $(L, \wedge, \top)$ is
  a commutative monoid and $a \wedge b \leq c$ iff $a \leq (b \Rightarrow c)$
  (residuation).  A \emph{distributive lattice} is a monoidal poset where
  $\wedge$ is the meet operation of a distributive lattice.
\end{definition}

Given a finite distributive lattice $L$, an $L$-\emph{enriched category}
$\cat{C}$ consists of a finite set of objects together with, for each pair
of objects $(x,y)$, an element $\Hom(x,y) \in L$ satisfying:
\begin{align*}
  \top &\leq \Hom(x,x) &&\text{(identity)} \\
  \Hom(y,z) \wedge \Hom(x,y) &\leq \Hom(x,z) &&\text{(composition)}
\end{align*}

The three lattices we use are:

\begin{enumerate}
\item The \textbf{two-element lattice} $\mathbf{2} = \{0, 1\}$ with
  meet $a \wedge b = \min(a,b)$ and unit $1$.
  A $\mathbf{2}$-enriched category records which hom-sets are nonempty.
\item The \textbf{cardinal lattice} $\mathbb{N}_\infty = \mathbb{N} \cup
  \{\infty\}$ with meet $a \wedge b = \min(a,b)$ and unit $\infty$.
  An $\mathbb{N}_\infty$-enriched category records the size of each
  hom-set.
\item The \textbf{power-set lattice} $\mathcal{P}(M)$ of a finite set $M$,
  with meet $A \wedge B = A \cap B$ and unit $M$.
  A $\mathcal{P}(M)$-enriched category with $|\Mor(\cat{C})| = |M|$
  records which specific morphisms belong to each hom-set, together with
  the composition table.
\end{enumerate}

The three-layer pipeline is:

\[
\cat{C}_1 \xmapsto{\text{refine}} \cat{C}_2 \xmapsto{\text{refine}} \cat{C}_3
\]

where $\cat{C}_1$ is $\mathbf{2}$-enriched (a \emph{support preorder}),
$\cat{C}_2$ is $\mathbb{N}_\infty$-enriched (a \emph{hom-count matrix}),
and $\cat{C}_3$ is a genuine finite category with a composition table.

\subsection{Enumerating support preorders}

A $\mathbf{2}$-enriched category with $k$ objects is precisely a
reflexive, transitive binary relation on $\{1,\dots,k\}$, i.e., a
preorder.  The number of labelled preorders on $k$ elements is given by
OEIS A000798: $1, 1, 4, 29, 355, 6942, \dots$.  We enumerate these via
the standard ``antichain extension'' method \cite{pfeiffer2004}: starting
from the empty poset, we add one element at a time, maintaining the poset
structure by choosing a lower set (downward-closed subset) and an upper
set (upward-closed subset) for the new element, subject to the constraints
that they are disjoint and every element of the lower set is below every
element of the upper set.

More precisely, given a relation $R$ on $\{0,\dots,q-1\}$ and a new element
$q$, we extend $R$ to $\{0,\dots,q\}$ by choosing:
\begin{itemize}
  \item $L \subseteq \{0,\dots,q-1\}$, a down-set of $R$;
  \item $U \subseteq \{0,\dots,q-1\}$, an up-set of $R$;
\end{itemize}
such that $L \cap U = \varnothing$ and $l < u$ in $R$ for all $l \in L$,
$u \in U$.  The extended relation $\tilde{R}$ is defined by:
\[
\tilde{R}(i,j) =
\begin{cases}
R(i,j) & i,j < q, \\
1 & i = j = q, \\
1 & j = q \text{ and } i \in L, \\
1 & i = q \text{ and } j \in U, \\
0 & \text{otherwise}.
\end{cases}
\]

The set of all extensions of $R$ is exactly the set of all preorders on
$\{0,\dots,q\}$ whose restriction to $\{0,\dots,q-1\}$ is $R$.

To obtain \emph{labelled} supports (where object indices matter), we first
generate canonical (unlabelled) support preorders via the antichain
extension of canonical poset quotients, then expand each into all $k!$
object-permuted versions, deduplicating by a word-sized bitmask.

\begin{proposition}
  \label{prop:canonical-supports}
  For $k \leq 8$, the set of labelled support preorders on $k$ objects
  can be generated in $O(k! \cdot C(k))$ time, where $C(k)$ is the number
  of canonical preorders on $k$ elements, by expanding each canonical
  preorder under all object permutations.
\end{proposition}

\subsection{Enumerating hom-count matrices}

Given a support preorder $R$, an $\mathbb{N}_\infty$-enriched category
refining $R$ is a $k \times k$ matrix $H$ of non-negative integers such
that $H(i,i) \geq 1$, $H(i,j) > 0 \iff R(i,j) = 1$, and
$\sum_{i,j} H(i,j) = n$, where $n$ is the total number of morphisms.
The composition axiom for $\mathbb{N}_\infty$:
\[
\min(H(j,\ell), H(i,j)) \leq H(i,\ell)
\]
is automatically satisfied for any matrix with non-negative entries and
$H(i,i) \geq 1$, since $\min(a,b) \leq \max(a,b) \leq H(i,\ell)$ only when
$H(i,\ell)$ is sufficiently large.  In practice this condition is almost
always satisfied and serves as a light filter.

We generate all matrices by distributing the $n - k$ ``extra'' morphisms
across the $k^2$ hom-cells via weak compositions.  The number of such
matrices for given $(n,k)$ is:
\[
\binom{n - k + k^2 - 1}{k^2 - 1} = \binom{n - k^2 - 1}{k^2 - 1}
\]

\subsection{Enumerating composition tables}

The third layer --- enriching over $\Set$ --- is the core combinatorial
search.  Given a hom-count matrix $H$, we fix $n$ morphisms and assign
each to a hom-cell according to $H$.  The composition table is then a
partial $n \times n$ matrix with entries in $\{0,\dots,n-1\} \cup \{-1\}$,
where $-1$ indicates a non-composable pair.  The search space is
constrained by:

\begin{itemize}
\item \textbf{Identities:} For each object $x$, the designated identity
  morphism $\id_x$ satisfies $\id_x ; f = f$ for all $f$ with
  $s(f) = x$, and $g ; \id_x = g$ for all $g$ with $t(g) = x$.
\item \textbf{Typing:} If $s(g) \neq t(f)$, then $f; g = -1$.
\item \textbf{Associativity:} For all composable triples $(f,g,h)$,
  $(f;g);h = f;(g;h)$.
\end{itemize}

We solve this constraint satisfaction problem via backtracking with
forward propagation.  At each step, we choose the composition cell with
the smallest domain (fewest candidate morphisms) and try each candidate
value, propagating associativity constraints to neighbouring cells.

\subsection{Cauchy completeness}

A category is \emph{Cauchy complete} (or \emph{idempotent-complete}) if
every idempotent endomorphism $e : x \to x$ (with $e;e = e$) splits:
there exist $r : x \to y$ and $s : y \to x$ such that $r;s = e$ and
$s;r = \id_y$.  Cauchy completeness is a natural finiteness condition
--- in a $\Set$-enriched setting, it corresponds to the Cauchy completeness
of the enriched category.

We enforce Cauchy completeness during search by a pruning predicate
$\texttt{partialCauchyViable}$: if a partially-filled composition table
already forces an idempotent (i.e., a cell $(e,e)$ has been assigned the
value $e$) and that idempotent has no potential splitting (no pair
$(r,s)$ of appropriate type that could witness the split), the search
branch is discarded.

For one-object categories, Cauchy completeness is equivalent to the
endomorphism monoid being a group.  We use this fact to construct
one-object Cauchy-complete categories directly via group-theoretic
constructions (Section~\ref{sec:decomposition}).

\section{Canonicalisation and quotient modes}
\label{sec:canonical}

Given a valid composition table, we produce a \emph{representative key}
that identifies the category up to a chosen equivalence relation.  The
key is a serialisation of the arrows-only data:
\[
\operatorname{key}(\cat{C}) = (k, n, \vec{s}, \vec{t}, \vec{\id}, \vec{;})
\]
where $\vec{s}$ and $\vec{t}$ are the source and target vectors,
$\vec{\id}$ is the identity vector, and $\vec{;}$ is the flat composition
table.

\subsection{Strict isomorphism}

Two categories $\cat{C}$ and $\cat{D}$ are \emph{strictly isomorphic}
if there exist bijections $\varphi : \Ob(\cat{C}) \to \Ob(\cat{D})$ and
$\psi : \Mor(\cat{C}) \to \Mor(\cat{D})$ preserving sources, targets,
identities, and composition.  To obtain a canonical key, we minimise the
key under all $k!$ object permutations and, for each object permutation,
all block-wise morphism permutations.  The inner morphism permutations
respect hom-block structure: for each pair of objects $(i,j)$, the
morphisms in $\Hom(i,j)$ can be permuted among themselves, with the
identity on each diagonal block fixed as the first element.

This is the Cruttwell--Leblanc exhaustive canonicalisation
\cite{cruttwell2014}, and for one-object categories it reduces to finding
a canonical generator order for the endomorphism group.

\subsection{Categorical equivalence}

Two categories $\cat{C}$ and $\cat{D}$ are \emph{categorically equivalent}
if there exists a functor $F : \cat{C} \to \cat{D}$ that is full, faithful,
and essentially surjective on objects.  Up to equivalence, every category
is equivalent to its \emph{skeleton}: a full subcategory containing exactly
one object from each isomorphism class.

The equivalence key of $\cat{C}$ is computed by first taking the skeleton
(identifying isomorphic objects via the existence of inverse pairs of
morphisms) and then applying the isomorphism key to the skeleton.

In disconnected categories, two objects in different connected components
may be structurally identical but lack connecting morphisms, which
prevents the skeleton computation from detecting their isomorphism.
This is the sole source of discrepancy between our equivalence counts and
published OEIS data; see Section~\ref{sec:results}.

\subsection{Labelled-object equality}

Under \emph{equality semantics}, object labels matter but morphism labels
within each hom-set do not.  Two categories are equal iff they have the
same object labels, hom-count matrix, and composition table.  This is the
finest quotient relation we consider; it corresponds to what would be
called ``labelled categories'' in the combinatorics literature.

The equality count for a category $\cat{C}$ with $k$ objects is
$k! / |\Aut(\cat{C})|$, where $\Aut(\cat{C})$ is the automorphism group
of $\cat{C}$ under object permutations.  The automorphism group size is
computed via object-class partitioning: two objects belong to the same
class iff they have the same hom-count profile, i.e., for all other
objects $t$:
\[
|\Hom(i,t)| = |\Hom(j,t)| \quad\text{and}\quad |\Hom(t,i)| = |\Hom(t,j)|
\]
If $H(i,i) = H(j,j)$ and the off-diagonal counts match, then $i$ and $j$
are placed in the same equivalence class.  The automorphism group is the
product of the symmetric groups on each class:
\[
|\Aut(\cat{C})| = \prod_{\text{classes } C} |C|!
\]
and the equality count is $k! / |\Aut(\cat{C})|$.

\subsection{Source--target duality}

A further quotient identifies each category $\cat{C}$ with its opposite
$\cat{C}^{\op}$, obtained by swapping sources and targets and reversing
composition order.  The dual key is:
\[
\operatorname{dualKey}(\cat{C}) = \min(\operatorname{key}(\cat{C}),
\operatorname{key}(\cat{C}^{\op}))
\]
The opposite of $\cat{C}$ is well-defined and self-dual categories are
those isomorphic to their opposite.

\section{Decomposition-based generation}
\label{sec:decomposition}

The shape-based search of Section~\ref{sec:enrichment} becomes
prohibitively slow for $n \geq 8$ because the number of hom-count
matrices grows as $\binom{n-k+k^2-1}{k^2-1}$ and each matrix's
composition-table search is exponential.  To scale beyond $n = 7$, we
decompose categories into smaller building blocks.

\subsection{Biconnected components}

A category is \emph{biconnected} if every hom-set is nonempty:
$\Hom(x,y) \neq \varnothing$ for all objects $x,y$.  The biconnected
components of a category are its connected components under the
equivalence relation generated by $\Hom(x,y) \neq \varnothing$.  Every
finite category factors as a ``gluing'' of its biconnected components
along profunctor edges.

More formally, let $\cat{C}$ be a finite category.  Define a relation
$\sim$ on $\Ob(\cat{C})$ by $x \sim y$ iff there exists a zigzag of
morphisms connecting $x$ and $y$ in both directions.  The equivalence
classes of $\sim$ are the \emph{support classes} of $\cat{C}$.  The
quotient poset $Q = \Ob(\cat{C}) / \sim$ inherits a partial order from
the support relation: $[x] \leq [y]$ if $\Hom(x,y) \neq \varnothing$ for
some representatives $x,y$.  For each class $[x]$, the full subcategory
on $[x]$ is biconnected.

\begin{proposition}
  \label{prop:decomposition}
  Every finite category $\cat{C}$ can be reconstructed from:
  \begin{enumerate}
  \item its biconnected components $\cat{B}_1, \dots, \cat{B}_q$,
  \item its quotient poset $Q$,
  \item and for each pair of classes $([i],[j])$ with $[i] \leq [j]$ in
    $Q$, a profunctor $\cat{B}_i^{\op} \times \cat{B}_j \to \Set$
    whose action determines the cross-component morphisms.
  \end{enumerate}
\end{proposition}

Our generation pipeline enumerates biconnected components via the
shape-based search (for $k$ up to $\lfloor\sqrt{n}\rfloor$, since a
biconnected category with $k$ objects has at least $k^2$ morphisms), then
assembles them into larger categories via skeleton enumeration
(Section~\ref{sec:skeletons}).

\subsection{One-object components: groups}

One-object Cauchy-complete categories are exactly finite groups.  We
construct these directly using the known classification of groups of
small orders, avoiding the generic composition-table search.  The
construction proceeds as follows:

\begin{enumerate}
\item \textbf{Abelian groups:} For each order $n$, factor $n$ into prime
  powers.  For each prime power $p^a$, enumerate integer partitions of
  $a$; each partition $\lambda_1 \geq \cdots \geq \lambda_r$ gives a group
  $\prod_i C_{p^{\lambda_i}}$.  Combine prime parts via direct products.
\item \textbf{Small non-abelian groups:} For orders $\leq 15$, use known
  constructors: dihedral groups $D_n$, dicyclic groups $Q_n$, and the
  alternating group $A_4$.
\item \textbf{P-groups:} For orders $p^3$, use the three abelian groups,
  the Heisenberg group of order $p^3$, and the nonabelian semidirect
  product $C_{p^2} \rtimes C_p$.
\item \textbf{Pq orders:} For orders $pq$ where $p|q-1$, construct
  semidirect products $C_q \rtimes C_p$ using the automorphism group of
  $C_q$.
\item \textbf{Fallback:} For unsupported orders, use generic
  Latin-square search with associativity constraints.
\end{enumerate}

\subsection{Skeleton assembly}
\label{sec:skeletons}

Given a quotient poset $Q$ on $q$ classes and a multiset of biconnected
components $\cat{B}_1, \dots, \cat{B}_q$, we construct a candidate
category $\cat{C}$ as follows:
\begin{enumerate}
\item Determine the hom-count matrix $M$ from the component sizes and
  poset relations: for $i \neq j$, if $[i] \leq [j]$ in $Q$, then
  $\Hom(\cat{B}_i, \cat{B}_j)$ receives at least one morphism;
  otherwise none.
\item Distribute any remaining morphisms (excess $n - \sum |\Mor(\cat{B}_i)|$)
  across the off-diagonal hom-cells.
\item Use the profunctor solver (Section~\ref{sec:profunctors}) to fill
  the composition data for cross-component morphisms.
\end{enumerate}

The skeleton canonical key ensures that only distinct isomorphism classes
are counted.  Crucially, when computing the canonical key, only classes
with identical component keys are permuted --- this prevents distinct
component-to-class assignments from colliding, a subtle bug that
previously caused significant undercounts at $(9,2)$.

\subsection{Shortcuts for small excess}

When the \emph{excess} $e = n - k$ is small ($e \leq 4$), we use a
fast path that constructs categories directly from groups and terminals,
avoiding the full skeleton assembly.  The shortcut enumerates all integer
partitions of $e$ and assigns each part to a group object (with order
equal to part $+ 1$), together with a parallel-arrow component.

For example, for $e = 1$:
\[
\{(2, C_2) \oplus 1\text{ terminal}\} \cup \{(3, \text{parallel arrows})\}
\]
For $e = 2$:
\[
\{(3, C_3)\} \cup \{(2, C_2) \oplus (2, C_2)\} \cup \{(4, \text{parallel arrows})\}
\]
For larger $e$, the partitions of $e$ generate all distributions of groups
across objects.

\section{The colimit--limit theorem}
\label{sec:theorem}

The decomposition of Section~\ref{sec:decomposition} is justified by a
general theorem about categories enriched over preorders.

Let $U : \Preord \to \Set$ be the forgetful functor.  For a preorder $P$,
write $U(P)\text{-Cat}$ for the category of categories enriched over the
monoidal poset $(U(P), \wedge, \top)$ where $\wedge$ is the meet and
$\top$ is the top element of $P$.

\begin{theorem}
  \label{thm:colimit-limit}
  The functor
  \[
  U(-)\text{-Cat} : \Preord^{\op} \longrightarrow \CAT/\Set
  \]
  sends all weighted colimits in the locally-posetal 2-category
  $\Preord$ to weighted limits in $\CAT/\Set$.
\end{theorem}

\begin{proof}[Proof sketch]
  For a fixed set $X$, let $\Preord(X)$ be the preorder of preorder
  relations on $X$ ordered by inclusion.  A $U(P)$-enriched category
  with object set exactly $X$ is the same thing as a monotone map
  $P \to \Preord(X)$.  Explicitly, $p \in P$ is sent to the preorder
  $x \leq_p y \iff p \in A(x,y)$ where $A(x,y)$ is the hom-object.

  The fibre over $X$ is therefore:
  \[
  (U(P)\text{-Cat})_X \cong \Preord(P, \Preord(X))
  \]
  If $L = \operatorname{colim}_j P_j$, then for every $X$:
  \[
  \Preord(L, \Preord(X)) \cong \lim_j \Preord(P_j, \Preord(X))
  \]
  Hence $U(L)\text{-Cat} \cong \lim_j U(P_j)\text{-Cat}$, not in
  plain $\Cat$ but in $\CAT/\Set$.
\end{proof}

\begin{corollary}
  \label{cor:limits}
  For a coproduct $P + Q$ of preorders:
  \[
  U(P+Q)\text{-Cat} \cong U(P)\text{-Cat} \times_{\Set} U(Q)\text{-Cat}
  \]
  In particular, the category of categories enriched over the disjoint
  union of two preorders is the pullback over $\Set$ of the categories
  enriched over each component.
\end{corollary}

This result justifies our decomposition pipeline: a category whose support
preorder is a disjoint union factors over its connected components, and
the factoring is captured by a pullback in $\CAT/\Set$.  Similarly,
gluing components along a profunctor corresponds to a pushout in
$\Preord$, dualising to a pullback in $\CAT/\Set$.

\section{Finite profunctors}
\label{sec:profunctors}

A \emph{finite profunctor} $P : \cat{A}^{\op} \times \cat{B} \to \Set$
between two finite categories $\cat{A}$ and $\cat{B}$ is a functor whose
fibres are finite sets.  Profunctors describe the cross-component
morphism structure in the decomposition pipeline.

We count profunctors of bounded total fibre size by solving a constraint
satisfaction problem over action variables:

\begin{itemize}
\item For each morphism $f : a' \to a$ in $\cat{A}$ and object $b$ in
  $\cat{B}$, a left action $P(f,b) : P(a,b) \to P(a',b)$.
\item For each object $a$ in $\cat{A}$ and morphism $g : b \to b'$ in
  $\cat{B}$, a right action $P(a,g) : P(a,b) \to P(a,b')$.
\end{itemize}

The actions must satisfy identity, composition, and the middle interchange
law:
\[
P(f, b') \circ P(a, g) = P(a', g) \circ P(f, b) : P(a,b) \to P(a',b')
\]

For one-object group profunctors, we use a faster algebraic approach:
the category of profunctors between groups $G$ and $H$ is equivalent to
the category of finite $(G,H)$-bisets.  We enumerate transitive bisets
via coset actions of $G \times H^{\op}$, then combine them as orbit
multisets.

\section{Implementation and performance}
\label{sec:implementation}

The algorithm is implemented in Haskell (GHC 9.6.7) with approximately
10,000 lines of code across 14 modules.  The key modules are:

\begin{itemize}
\item \texttt{CodexSlop.Category} --- data type, validation, serialisation.
\item \texttt{CodexSlop.Shape} --- hom-count matrices, support preorders.
\item \texttt{CodexSlop.Search} --- backtracking composition-table solver.
\item \texttt{CodexSlop.Canonical} --- key generation for all quotient modes.
\item \texttt{CodexSlop.Generate} --- decomposition-based generation.
\item \texttt{CodexSlop.Enrichment} --- lattice typeclass and enrichment
  enumeration.
\item \texttt{CodexSlop.Group} --- group construction.
\item \texttt{CodexSlop.Profunctor} --- profunctor counting.
\end{itemize}

Caching is used at multiple levels: biconnected component keys
(``.cat-enum-cache/v1/biconnected/''), generated isomorphism keys
(``v4/generated/isomorphism/''), and profunctor profile counts
(``v1/profunctor-counts/'').  First runs are significantly slower
(populating the cache), while subsequent runs are near-instantaneous.

\subsection{Verified counts}
\label{sec:results}

Table~\ref{tab:counts} shows the number of Cauchy-complete finite
categories up to strict isomorphism for $n \leq 12$.

\begin{table}[h]
\centering
\begin{tabular}{c|c|c|l}
$n$ & Total & OEIS match & Notes \\
\hline
1 & 1 & --- & \\
2 & 2 & --- & \\
3 & 4 & --- & \\
4 & 11 & --- & \\
5 & 25 & --- & \\
6 & 63 & --- & \\
7 & 163 & --- & \\
8 & 451 & --- & \\
9 & 1312 & 2+278+458+371+151+39+10+2+1 & See below \\
10 & 3359 & --- & \\
11 & 9309 & --- & \\
12 & 20341 & --- &
\end{tabular}
\caption{Cauchy-complete categories up to strict isomorphism}
\label{tab:counts}
\end{table}

The OEIS row for $n=9$ (A384134) lists $2, 278, 457, 371, 151, 39, 10,
2, 1$ for a total of 1311.  Our count of 458 at $k=3$ (vs.~OEIS's 457)
is the result of a genuine mathematical edge case: two categories with
hom-count matrices

\[
\begin{pmatrix}
1 & 1 & 2 \\
1 & 1 & 2 \\
0 & 0 & 1
\end{pmatrix}
\quad\text{and}\quad
\begin{pmatrix}
1 & 2 & 2 \\
0 & 1 & 1 \\
0 & 1 & 1
\end{pmatrix}
\]

that are formal duals (each is isomorphic to the opposite of the other)
but not isomorphic (they have different object-out/in multisets:
$\{(4,2),(4,2),(1,5)\}$ vs $\{(5,1),(2,4),(2,4)\}$).  Under strict
isomorphism, both should count separately; under source--target duality,
they collapse to one.

Table~\ref{tab:equivalence} shows the comparison between strict
isomorphism, categorical equivalence, and isomorphism-with-duality
for $n=9$.

\begin{table}[h]
\centering
\begin{tabular}{c|r|r|r}
$k$ & Strict isomorphism & Categorical equivalence & Isomorphism + dual \\
\hline
1 & 2 & 2 & 2 \\
2 & 278 & 278 & 186 \\
3 & 458 & 457 & 264 \\
4 & 371 & 370 & 217 \\
5 & 151 & 150 & 95 \\
6 & 39 & 39 & 39 \\
7 & 10 & 10 & 8 \\
8 & 2 & 2 & 2 \\
9 & 1 & 1 & 1 \\
\hline
\text{Total} & 1312 & 1309 & 870
\end{tabular}
\caption{Comparison of quotient modes for $n=9$}
\label{tab:equivalence}
\end{table}

\subsection{Theoretical complexity}

The shape-based search has worst-case complexity
$O(\binom{n-k+k^2-1}{k^2-1} \cdot d^{n^2})$ where $d$ is the average
hom-block size.  The decomposition-based generation reduces this by
exploiting the biconnected-component factorisation: the search is
confined to components with $k \leq \lfloor\sqrt{n}\rfloor$, and the
assembly step is polynomial in the number of components.

The equality count via object-class partitioning is $O(k^2 \cdot n)$ per
category, compared to $O(k! \cdot n^2)$ for the naive hash-based approach
--- an improvement of factorial-to-polynomial.

\section{Conclusion and future work}
\label{sec:conclusion}

We have presented a complete algorithmic framework for the exact
enumeration of finite categories, with verified counts up to 12
morphisms.  The three-layer stratification, decomposition-based
generation, and multiple quotient modes provide a flexible tool for
exploring the structure of finite categories.

Several directions for future work remain:

\begin{enumerate}
\item \textbf{Higher counts.}  Extending the enumeration to $n > 12$
  requires optimising the biconnected component search for large excess
  values.  The current bottleneck is the shape-based enumeration for
  two-object biconnected categories with $n \geq 13$.
\item \textbf{Profunctor decomposition.}  The colimit--limit theorem
  (Theorem~\ref{thm:colimit-limit}) suggests a systematic approach to
  decomposing categories over their support preorders.  Making this
  decomposition algorithmic would further reduce the search space.
\item \textbf{Distributive-lattice enrichment.}  The enrichment framework
  (Section~\ref{sec:enrichment}) allows categories enriched over
  arbitrary finite distributive lattices.  Exploring these spaces may
  reveal connections to other combinatorial structures.
\item \textbf{Column-regular categories.}  The discrepancy between our
  equivalence counts and OEIS data at $n=9$, $k \geq 4$ arises from
  disconnected categories whose isomorphic components lack connecting
  morphisms.  A proper treatment of this case would compute equivalence
  by comparing component structures rather than requiring explicit
  isomorphisms.
\end{enumerate}

\begin{thebibliography}{10}

\bibitem{baez2001}
J.~Baez and J.~Dolan, ``Categorification,'' in \emph{Higher category
  theory}, Contemp. Math.~230, 2001.

\bibitem{besche2002}
H.~U. Besche, B.~Eick, and E.~A. O'Brien, ``A millennium project:
  constructing small groups,'' \emph{Int. J. Algebra Comput.}, 2002.

\bibitem{cruttwell2014}
G.~Cruttwell, ``Counting finite categories,'' \emph{AMS Spring Eastern
  Meeting}, 2014.

\bibitem{distler2011}
A.~Distler and J.~D. Mitchell, ``The number of nilpotent semigroups of
  degree $3$,'' \emph{Electron. J. Combin.}, 2011.

\bibitem{jevons1897}
W.~S. Jevons, ``On the principle of the classification of finite
  distributive lattices,'' \emph{Proc. London Math. Soc.}, 1897.

\bibitem{pfeiffer2004}
G.~Pfeiffer, ``Counting transitive relations,'' \emph{J. Integer Seq.},
  2004.

\bibitem{smallcats}
``SmallCategories,'' \url{https://smallcats.info/about}.

\end{thebibliography}

\end{document}
